\begin{table}[htbp]\centering
\begin{threeparttable}
\caption{Independent cultures: the rules as fixed, on 52 LongGloDB cultures from other collections.}
\label{tab:lgdb}
\footnotesize\setlength{\tabcolsep}{4pt}
\begin{tabular}{L{0.11\linewidth}L{0.14\linewidth}L{0.09\linewidth}L{0.06\linewidth}L{0.19\linewidth}L{0.13\linewidth}L{0.10\linewidth}}
\toprule
Analysis & LSU reference & False known (of 102) & Share & Absent species: genus right/wrong/none (of 59) & Absent genera: genus given (of 43) & Known kept (of 164)\\
\midrule
Main & intact & 4 (4) & 3.9\% & 6 / 0 / 53 & 0 & 136 (151)\\
 & removed & 4 (4) & 3.9\% & 6 / 0 / 53 & 0 & --\\
Second & intact & 5 (5) & 4.9\% & 29 / 0 / 30 & 18 & 139 (154)\\
 & removed & 5 (5) & 4.9\% & 29 / 0 / 30 & 18 & --\\
\bottomrule
\end{tabular}
\begin{tablenotes}\footnotesize
\item With the ITS distance rule; in brackets, the same without it. \emph{LSU reference}: intact, or with each absent
species removed in its own round. \emph{False known}: ASVs of species the reference
does not hold called a known species or species group. \emph{Absent species}: ASVs of species the reference does not
hold, in genera it holds. \emph{Absent genera}: ASVs of genera it does not hold, so any genus given is wrong.
\emph{Known kept}: ASVs of species the reference holds whose call contains their species; LSU reference intact
only. \emph{Second}: the second analysis, fixed in advance, with VT coverage counted over the part of each MaarjAM
reference that the ASV spans.
\end{tablenotes}
\end{threeparttable}
\end{table}
